\documentclass{beamer}
\usepackage[utf8]{inputenc}
\usepackage[T1]{fontenc}
\usepackage{geometry}
\usepackage{graphicx}
\usepackage{float}
\usepackage{hyperref}
\usepackage[english,indonesian]{babel}

\title{Optimasi Pemrosesan Template LaTeX Menggunakan Pendekatan Berbasis AI}
\author{Academic Support Agent}
\date{\today}

\begin{document}
\maketitle

\begin{abstract}

\end{abstract}

\chapter{PENDAHULUAN}


\chapter{TINJAUAN PUSTAKA}


% Original content (commented out):
% % !TeX spellcheck = en_B
% !TeX encoding = UTF-8 %%%These comments just make you look like you know what you're doing
\documentclass[8pt]{beamer}  %%% Specifies the class, with 8pt font size

\input{preamble} %%%adds the general preamble

\title{{\color{aa}\Huge\adfbullet{9}}AL FM Discrete} %%%Title and subtitle
\subtitle{Template, \textattachfile{Template.tex}{(TeX)}} %%%This embeds the tikz source code into the resulting pdf...

\begin{document}

\frame{\titlepage} %%%Makes titlepage

\setlength{\abovedisplayskip}{0pt}
\setlength{\belowdisplayskip}{0pt}
\setlength{\abovedisplayshortskip}{0pt}
\setlength{\belowdisplayshortskip}{0pt}  %%%Compresses math

\begin{frame}{Title}   %%%Frame with title
    Here is some normal text. Followed by a list.
\begin{itemize}
	\item First item.
	\item Second item
    \item \begin{itemize}
            \item<+-> List inside a list!!
            \item<+-> And another. These appear one at a time.
        \end{itemize} 
    \item Final item
\end{itemize}

\alert<+->{Here is an alert which is only highlighted on the last slide}

\end{frame}

\begin{frame}{Boxes}
	There are three commonly used boxes.
	
	\begin{definition}
	This is a \textbf{definition box}, it can be helpful to put the word being defined into bold text.
	\end{definition}
	
	\begin{problem}
	        This is a problem box. With some displayed math...
	        \[y=x^2\]
	\end{problem}
	
	\begin{solution}<2->
	This is a solution which only appears on the next slide.
	\end{solution}
	
	\begin{solution}<3->
	Appears even later!
	\end{solution}

\end{frame}

\begin{frame}{The sol command}

The first time you use \sol{it doesn't cover it up}.

However the second time \sol{it does}.

You can put math inside it, $5x+4x=$ \sol{$9x$}

And you can use it inside a math environment:

\begin{align*}
    y &= \sol{5} \\
    x &= \sol{$x^2$} \\
    xy &= \sol{$5x^2$}
\end{align*}

But this will throw up some errors, but still compile.

\end{frame}

\begin{frame}{columns}
    \begin{columns}
    \begin{column}{.5\linewidth}
    Columns can be very useful.
    \end{column}
    \begin{column}{.5\linewidth}
    To display info side by side.
    \end{column}
    \end{columns}
        \begin{columns}[T] %%%Aligned at the top!!!
    \begin{column}{.5\linewidth}
    \begin{problem}
            You can even put boxes inside the columns to save space.
    \end{problem}
    \end{column}
    \begin{column}{.5\linewidth}
    \begin{solution}<2->
    With a solution.
    \end{solution}
    \end{column}
    \end{columns}
    
\begin{definition}
\noindent
\begin{minipage}{.4\linewidth}
Using columns inside boxes is best achieved with minipages.
\end{minipage}%
\begin{minipage}{.6\linewidth}
\begin{center}
\colorbox{cc!30}{
\begin{nicetable}{cc|ccc}
\multicolumn{2}{c}{} & \multicolumn{3}{c}{Player $2$}\\
\multicolumn{1}{c}{} &  & $X$  & $Y$ & $Z$ \\ \cline{2-5}
\raisebox{0.0cm}{\multirow{3}*{\rotatebox{90}{Player $1$}}}  & $P$ & $4$ & $2$ & $2$ \\
& $Q$ & $-3$ & $5$ & $1$ \\
& $R$ & $2$ & $-1$ & $3$ \\
\end{nicetable}}
\end{center}
\end{minipage}
\end{definition}
\end{frame}







\begin{frame}{The tikzmarknode command}

    \alert{Useful when annotating equation}
    
    Use it anywhere to create a node named anything that you can refer back to later. Also compiles but with errors...
    
	\begin{flalign*}
		\text{Maximise} && P &= V-3 && \\
		\text{Subject to} && V -6p-4q + s_1 &= 0 && \\
				  && V -p-7q + s_2 &= 0 && \\
				  && \tikzmarknode{A}{p + q + s_3} &= 1 &&
	.\end{flalign*}

\begin{tikzpicture}[overlay,remember picture]
    \draw[color=aa,<-] (A) --++ (1,-1) node[fill=cc] {The annotation $x^2$};
\end{tikzpicture}

\end{frame}

\begin{frame}{Tables}

\begin{center}
\colorbox{cc!30}{
\begin{nicetable}{c|cccccc|c}
$P$ & $v$ & $p$ & $q$ & $s_1$ & $s_2$ & $s_3$ & RHS \\ 
  \hline
1 & $-1$ & $0$ & $0$ & 0 & 0 & 0 & $-3$   \\ 
   \hline
0 & 1 & $-6$ & $-4$ & 1 & 0 & 0 & 0 \\ 
  0 & 1 & $-1$ & $-7$ & 0 & 1 & 0 & 1  \\ 
  0 & 0 & $1$ & 1 & 0 & 0 & 1 & 1 \\ 
\end{nicetable}}
\end{center}

Tikz pictures, as well as tables, basically everything can be shrunk with the adjust box command. You can specify max width, max height, width height etc.

\centering
\adjustbox{max width=4cm}{
\colorbox{cc!30}{
\begin{nicetable}{c|cccccc|c}
$P$ & $v$ & $p$ & $q$ & $s_1$ & $s_2$ & $s_3$ & RHS \\ 
  \hline
1 &\boxed{$-1$} & $0$ & $0$ & 0 & 0 & 0 & $-3$   \\ 
   \hline
0 & 1 & $-6$ & $-4$ & 1 & 0 & 0 & 0 \\ 
  0 & 1 & $-1$ & $-7$ & 0 & 1 & 0 & 0  \\ 
  0 & 0 & $1$ & 1 & 0 & 0 & 1 & 1 \\ 
\end{nicetable}}
}

\end{frame}


\begin{frame}[shrink=60]{The shrink command, automatically shrinks to fit on page, but is helpful to give a number so it knows how much horizontal space it has.}
This is the contents.

\lipsum

\begin{problem}
        \lipsum
\end{problem}

\end{frame}


\begin{frame}{The onslide}
    Can be used on anything but really cool with tikz pictures.
                \centering
\adjustbox{max width=.8\textwidth}{
\def\complex{(0,0) ellipse (7cm and 6cm)} %%%%These can be used as variables
\def\real{(0,0) ellipse (6.5cm and 3.5cm)}
\def\algebraic{(1.7,0) circle (5cm)}
\def\imag{(1.3,-4.7) ellipse (2.2cm and 1cm)}
\def\int{(2.9,-0.6) ellipse (2.6cm and 2.3cm)}
\def\whole{(2.9,-0.7) ellipse (2.1cm and 1.6cm)}
\def\count{(2.9,-0.9) ellipse (1.3cm and 1cm)}
\begin{tikzpicture}
        %%Complex%%%%%%%%%%%%%%%%%%5
        \onslide<2->{\filldraw[orange,fill opacity=0.1] \complex;
        \node[color=orange] at (0,5.5) {Complex Numbers};
        \node[color=orange] at (-4,4.2) {$\pi+3i$};
        \node[color=orange] at (-5,3.2) {$e-\pi i$};}
        %%%%%%Real%%%%%%%%%%%%%5
        \onslide<3->{\filldraw[green,fill opacity=0.1] \real;
        \node[color=green] at (0,3) {Real Numbers};}
        %%%%Transcendental%%%%%%%%%%5
        \onslide<4->{\node[color=green,align=left] at (-5,0.3) {Transcendental \\ Numbers};
        \node[color=green] at (-6,-0.5) {$\pi$};
        \node[color=green] at (-5.3,-1) {$e$};}      
        %%%%%%%Irrational and real%%%%%%
        \onslide<5->{\node[color=green] at (-2.9,2.3) {Irrational Numbers};
        \node[color=green] at (2.9,2.3) {Rational Numbers};
        \draw[color=green] (0,-3.5) -- (0,2.5);
        \node[color=green] at (-1.7,1.4) {$\sqrt{2}$};
        \node[color=green] at (-2.6,-0.3) {$-\sqrt{5}$};                
        \node[color=green] at (-1.3,-2.2) {$\sqrt[3]{26}$};
        \node[color=green] at (0.9,1.6) {$\frac{2}{3}$};                
        \node[color=green] at (5.7,0.7) {$\frac{97}{32}$};}
        %%%%%Algebraic%%%%%%%%%%
        \onslide<6->{\filldraw[blue,fill opacity=0.1]  \algebraic;
        \node[color=blue] at (1.7,4.5) {Algebraic Numbers};
        \node[color=blue] at (3.7,3.9) {$9+5i$};
        \node[color=blue] at (4.6,3.1) {$1+i$};}
        %%%%%Imag%%%%%%%%%%%%%%%
        \onslide<7->{\filldraw[purple,fill opacity=0.1]  \imag;
        \node[color=purple] at (1.3,-4.3) {\footnotesize Pure Imaginary Numbers};
        \node[color=purple] at (-0.5,-5) {$e i$};
        \node[color=purple] at (0.5,-5.3) {$\pi i$};
        \node[color=purple] at (0.8,-4.6) {$i$};
        \node[color=purple] at (2.0,-4.7) {$5 i$};}
        %%%%%Int%%%%%%%%%%%%%
        \onslide<8->{\filldraw[color=blue, fill opacity=0.1] \int;
        \node[color=blue] at (2.9,1.2) {Integers};
        \node[color=blue] at (0.9,0.4) {$5$};
        \node[color=blue] at (2.6,-2.6) {$-99$};}
        %%%%%5Whole%%%%%%5
        \onslide<9->{\filldraw[color=blue, fill opacity=0.1] \whole;
        \node[color=blue] at (2.9,0.3) {Whole Numbers};
        \node[color=blue] at (4.5,-0.6) {$0$};}
        %%%%%%%%%%Count%%%%%%%%%
        \onslide<10->{\filldraw[color=blue, fill opacity=0.1] \count;
        \node[color=blue,align=center] at (2.9,-0.5) {Counting \\ Numbers};
        \node[color=blue] at (2.2,-1.2) {$1$};
        \node[color=blue] at (2.8,-1.6) {$2$};
        \node[color=blue] at (3.4,-1.3) {$99$};}
\end{tikzpicture}}

\end{frame}

\end{document}


\end{document}